% answering.tex -- Module 8: Answering HSC questions (NESA key verbs + model answers).
% Verb meanings are written in our own simple words (not the NESA glossary text); they are
% the same as in the other books of the series (mod5, mod6).
% Model answers: compare (Ch 1/3), explain (Ch 2), analyse (Ch 4), evaluate (Ch 5),
% calculate (Ch 4).
%   Thomson: e/m_e = 1.759e11 C/kg, 1837 x the hydrogen ion's q/m (9.576e7 C/kg).
%   H-alpha (3 -> 2) from the Rydberg equation: 656.3 nm; hcR = 13.61 eV (assess row).
%   Binding energy (values from ch4.tex): fission 1786 -> 1970 MeV (+184 MeV); D-T 17.6 MeV;
%   per nucleon released: fission 0.78 MeV, D-T 3.5 MeV.
%   Tc-99m: lambda = 0.1155 h^-1; ln(800/300) = 0.9808; t = 8.490 h = 8 h 29 min.
\unnumberedchapter[\faPenNib]{Answering HSC questions}

% marking-point tag used in the model answers below
\newcommand\ansPt[1]{\textcolor{hscbrand}{\bfseries(#1)}}

Every HSC question starts with a \textbf{key verb}, such as ``explain'' or ``evaluate''.
The verb tells you what kind of answer the marker wants (see the table below). Then count
the marks: give about \textbf{one correct point for each mark}. Write one idea in each
sentence, and link the ideas with ``so'' or ``because''. Many Module~8 questions are about
experiments: say what was \emph{done}, what was \emph{observed} and what it \emph{showed}.
In the model answers below, each point that earns a mark is numbered, like \ansPt{1}.

\begin{table}
  % no numbered caption: a number here would print as ``Table 5.x'' (chapter counter 5)
  \caption*{The key verbs used in HSC Physics questions}
  \centering
  \small
  \begin{tabular}{@{}L{18mm}L{63mm}L{60mm}@{}}
    \toprule
    \colhead{Verb} & \colhead{What it asks you to do} & \colhead{Module 8 example}\\
    \midrule
    \textbf{Identify} & Name the thing. A word or a short phrase is enough.
      & Identify the exchange particle of the strong force. \emph{(The gluon.)}\\
    \textbf{Outline} & Give the main points only, in order. No details.
      & Outline the proton--proton chain. \emph{(Protons fuse to make deuterium, then
        helium-3, then helium-4.)}\\
    \textbf{Describe} & Say what it is like or what happens. You do not need to say why.
      & Describe the cosmic microwave background. \emph{(Faint microwaves from every
        direction, like a black body at \qty{2.7}{K}.)}\\
    \textbf{Explain} & Say \emph{why} or \emph{how} it happens. Link each cause to its
      effect with ``because'' or ``so''.
      & Model answer~2 below.\\
    \textbf{Compare} & Give the similarities \emph{and} the differences.
      & Model answer~1 below.\\
    \textbf{Analyse} & Break it into parts (steps, particles, numbers). Show how they
      are linked and what follows.
      & Model answer~3 below.\\
    \textbf{Assess} & Make a judgement (is it true? how good? how big?) and give reasons.
      & Assess the Bohr model. \emph{(It explains hydrogen's lines, but it fails for
        other atoms. Useful, but limited.)}\\
    \textbf{Evaluate} & Judge it against criteria: give strengths \emph{and} limitations,
      then a clear final judgement.
      & Model answer~4 below.\\
    \textbf{Justify} & Give reasons or evidence that support an answer or a decision.
      & Justify using a moderator in a reactor. \emph{(It slows the neutrons, so U-235
        captures them and the chain reaction goes on.)}\\
    \textbf{Calculate} & Work out a number. Show the formula, the substitution with units,
      and the answer with a unit.
      & Model answer~5 below.\\
    \bottomrule
  \end{tabular}
\end{table}

\section*{Model answer 1: compare (3 marks)}

\emph{Compare the emission spectrum and the absorption spectrum of hydrogen gas.}

\begin{modelanswer}
\ansPt{1} \textbf{Similarity:} both are line spectra, with lines at the \emph{same}
wavelengths (for example \qty{656}{nm}). Both come from jumps between the same energy
levels: $hf = E_\mathrm{i} - E_\mathrm{f}$.

\ansPt{2} \textbf{Emission:} a hot gas at low pressure gives bright lines on a dark
background. Its electrons drop to lower levels and give out photons.

\ansPt{3} \textbf{Absorption:} white light through a cooler gas gives dark lines on a
continuous spectrum. The atoms absorb only photons that match a gap, and give them out
again in all directions.
\end{modelanswer}

\begin{markingcriteria}
  \criterion{A similarity and two differences, each linked to energy levels}{3}
  \criterion{A similarity and one difference, or two differences}{2}
  \criterion{One correct feature of either spectrum}{1}
\end{markingcriteria}

\section*{Model answer 2: explain (3 marks)}

\emph{Explain how Thomson's experiment showed that cathode rays are particles that are much
lighter than atoms.}

\begin{modelanswer}
\ansPt{1} Electric and magnetic fields bend cathode rays, but not light. So the rays are
charged particles. They bend towards the positive plate, so they are negative.

\ansPt{2} Thomson set crossed fields so that the beam went straight: $qE = qvB$, so
$v = \dfrac{E}{B}$. The magnetic field alone bent the beam into a circle:
$qvB = \dfrac{mv^2}{r}$, so $\dfrac{q}{m} = \dfrac{E}{B^2 r}$.

\ansPt{3} This $q/m$ is about 1800 times that of a hydrogen ion. With the same size
of charge, the particle has about $\tfrac{1}{1800}$ of the mass of the lightest atom. Every
cathode metal gave the same value, so these electrons are in all atoms.
\end{modelanswer}

\begin{markingcriteria}
  \criterion{Links the deflection, the method for $q/m$ and the comparison with hydrogen
    (points 1--3)}{3}
  \criterion{Two of these points, linked}{2}
  \criterion{One relevant point, for example ``the rays are bent by fields''}{1}
\end{markingcriteria}

\section*{Model answer 3: analyse (4 marks)}

\emph{Use the binding energy per nucleon curve to analyse why both the fission of
uranium-235 and the fusion of hydrogen nuclei release energy.}

\begin{modelanswer}
\ansPt{1} Binding energy per nucleon shows how tightly a nucleus is held. The curve
(\cref{ch4:fig:becurve}) rises steeply to helium-4 (\qty{7.1}{MeV}) and peaks at iron-56
(\qty{8.8}{MeV}). Then it falls slowly to uranium (\qty{7.6}{MeV}).

\ansPt{2} \textbf{Fission:} uranium-235 (\qty{7.6}{MeV} per nucleon) splits into two
middle-sized nuclei (8.3 to \qty{8.6}{MeV} per nucleon). The total binding energy rises by
about \qty{184}{MeV}.

\ansPt{3} \textbf{Fusion:} hydrogen-2 (\qty{1.1}{MeV} per nucleon) and hydrogen-3 join to
make helium-4 (\qty{7.1}{MeV} per nucleon). About \qty{17.6}{MeV} is released.

\ansPt{4} \textbf{Link:} both reactions move the products up the curve, towards iron.
More tightly bound products have less mass. The lost mass becomes kinetic energy:
$E = \Delta m\,c^2$. The curve is steepest for light nuclei, so fusion releases more energy
per nucleon (\qty{3.5}{MeV}) than fission (\qty{0.8}{MeV}).
\end{modelanswer}

\begin{markingcriteria}
  \criterion{Curve, both reactions with values, and the link to mass lost (points 1--4)}{4}
  \criterion{Both reactions and the link to mass, but few values}{3}
  \criterion{Applies the curve to one reaction, or states the link to mass only}{2}
  \criterion{One relevant point, for example ``iron is at the peak''}{1}
\end{markingcriteria}

\section*{Model answer 4: evaluate (4 marks)}

\emph{``The Standard Model is a complete theory of matter.'' Evaluate this claim.}

\begin{modelanswer}
\ansPt{1} \textbf{Strengths:} it has 6 quarks, 6 leptons and the exchange bosons of three
forces. The quarks explain every hadron (baryons $qqq$, mesons $q\overline{q}$). The
W boson explains beta decay: a down quark becomes an up quark and gives out a
$\mathrm{W}^-$.

\ansPt{2} \textbf{Evidence:} accelerators found the particles it predicted: the W and Z
bosons (CERN, 1983), the top quark (Fermilab, 1995) and the Higgs boson (LHC, 2012).
\tcbbreak % layout QA (book build): break before point 3, so "(continued)" is not two lone lines

\ansPt{3} \textbf{Limitations:} it leaves out gravity. It does not explain neutrino mass,
dark matter or dark energy (about \qty{95}{\percent} of the Universe). It cannot say why
there is far more matter than antimatter.

\ansPt{4} \textbf{Judgement:} the model is very successful for what it covers, but it is
incomplete. So the claim is \textbf{false}: our understanding of matter is still
incomplete.
\end{modelanswer}

\begin{markingcriteria}
  \criterion{Strengths and limitations with evidence, and a clear judgement (points 1--4)}{4}
  \criterion{Strengths and limitations with a judgement, but little evidence}{3}
  \criterion{A strength or a limitation of the model}{2}
  \criterion{One relevant point about the Standard Model}{1}
\end{markingcriteria}
% layout QA (book build): one extra line on this page keeps Model answer 5 and its
% marking criteria together on the page before "Formula sheet explained".
\enlargethispage{\baselineskip}

\section*{Model answer 5: calculate (3 marks)}

\emph{At 8:00 am a hospital prepares a dose of technetium-99m with an activity of
\qty{800}{MBq}. Its half-life is \qty{6.0}{h}. A scan needs an activity of at least
\qty{300}{MBq}. Calculate the latest time at which the scan can start.}

\begin{modelanswer}
\ansPt{1} Find the decay constant (keep the time in hours):
$\lambda = \dfrac{\ln 2}{t_{\frac{1}{2}}} = \dfrac{0.6931}{\qty{6.0}{h}}
= \qty{0.1155}{h^{-1}}$.

\ansPt{2} Activity follows the decay law, $A_\mathrm{t} = A_0 e^{-\lambda t}$. Divide by
$A_0$, then take $\ln$ of both sides:
\begin{align*}
  e^{-\lambda t} &= \frac{\qty{300}{MBq}}{\qty{800}{MBq}} = 0.375 && \text{divide by $A_0$}\\
  -\lambda t &= \ln 0.375 = -0.9808 && \text{take $\ln$ of both sides}\\
  t &= \frac{0.9808}{\qty{0.1155}{h^{-1}}} = \qty{8.49}{h} && \text{divide by $-\lambda$}
\end{align*}

\ansPt{3} \final{$t = \qty{8.5}{h}$}: \qty{8.49}{h} is 8~h 29~min, so the scan must
start by about \textbf{4:29 pm}. Check: \qty{8.5}{h} is between one half-life
(\qty{400}{MBq}) and two half-lives (\qty{200}{MBq}) \cmark.
\end{modelanswer}

\begin{markingcriteria}
  \criterion{Correct decay constant, with a time unit that matches $t$}{1}
  \criterion{Correct substitution and rearranging with $\ln$}{1}
  \criterion{Correct time, with a unit (or as a clock time)}{1}
\end{markingcriteria}
